\documentclass{article}
\usepackage[utf8]{inputenc}
\usepackage{hyperref}
\usepackage{listings}
\usepackage{xcolor}

\definecolor{codegreen}{rgb}{0,0.6,0}
\definecolor{codegray}{rgb}{0.5,0.5,0.5}
\definecolor{backcolour}{rgb}{0.95,0.95,0.92}

\lstdefinestyle{mystyle}{
    backgroundcolor=\color{backcolour},
    commentstyle=\color{codegreen},
    keywordstyle=\color{blue},
    numberstyle=\tiny\color{codegray},
    stringstyle=\color{codegreen},
    basicstyle=\ttfamily\footnotesize,
    breakatwhitespace=false,
    breaklines=true,
    captionpos=b,
    keepspaces=true,
    numbers=left,
    numbersep=5pt,
    showspaces=false,
    showstringspaces=false,
    showtabs=false,
    tabsize=2
}

\lstset{style=mystyle}

\title{FFXIV Character Search - User Guide}
\author{Development Team}
\date{\today}

\begin{document}

\maketitle

\section{Introduction}
FFXIV Character Search is a React-based web application that allows users to search for and view character information from Final Fantasy XIV using the Lodestone website.

\section{Prerequisites}
Before running the application, ensure you have the following installed:
\begin{itemize}
    \item Node.js (version 14 or higher)
    \item npm (usually comes with Node.js)
    \item Git (for version control)
\end{itemize}

\section{Local Development}

\subsection{Clone the Repository}
First, clone the repository using git:

\begin{lstlisting}[language=bash]
git clone <repository-url>
cd ffxiv-character-search
\end{lstlisting}

\subsection{Install Dependencies}
Install all required dependencies using npm:

\begin{lstlisting}[language=bash]
npm install
\end{lstlisting}

\subsection{Running the Development Server}
To start the development server:

\begin{lstlisting}[language=bash]
npm run dev
\end{lstlisting}

This will:
\begin{itemize}
    \item Start the development server on port 3000
    \item Open your default browser to http://localhost:3000
    \item Enable hot reloading for development
\end{itemize}

\section{Building for Production}
To create a production build:

\begin{lstlisting}[language=bash]
npm run build
\end{lstlisting}

This will create optimized files in the \texttt{public/dist} directory.

\section{GitHub Pages Deployment}

\subsection{Repository Setup}
1. Create a new repository on GitHub
2. Push your code:
\begin{lstlisting}[language=bash]
git init
git add .
git commit -m "Initial commit"
git branch -M main
git remote add origin https://github.com/username/ffxiv-character-search.git
git push -u origin main
\end{lstlisting}

\subsection{Enable GitHub Pages}
1. Go to your repository's Settings
2. Navigate to Pages section
3. Select:
   \begin{itemize}
       \item Source: Deploy from a branch
       \item Branch: main
       \item Folder: /public
   \end{itemize}
4. Click Save

Your site will be available at:\\
\url{https://username.github.io/ffxiv-character-search}

\section{Using the Application}

\subsection{Searching for Characters}
1. Enter a character name (required)
2. Optionally enter:
   \begin{itemize}
       \item World name
       \item Data Center name
   \end{itemize}
3. Click Search
4. Results will display below the search form

\subsection{Viewing Character Details}
1. From search results, click "Select Character"
2. The character's details will display, including:
   \begin{itemize}
       \item Character portrait and basic information
       \item Combat jobs (Tanks, Healers, DPS)
       \item Crafting and Gathering jobs
   \end{itemize}

\section{Project Structure}
\begin{lstlisting}
ffxiv-character-search/
├── public/
│   ├── css/
│   ├── dist/
│   └── index.html
├── src/
│   ├── components/
│   │   ├── CharacterSearch.jsx
│   │   ├── SearchForm.jsx
│   │   ├── SearchResults.jsx
│   │   └── CharacterDetails.jsx
│   ├── utils/
│   │   └── api.js
│   └── App.jsx
├── package.json
└── webpack.config.js
\end{lstlisting}

\section{Troubleshooting}

\subsection{Common Issues}
\begin{itemize}
    \item If the development server fails to start:
        \begin{itemize}
            \item Check if port 3000 is already in use
            \item Try \texttt{npm install} again
        \end{itemize}
    \item If character searches fail:
        \begin{itemize}
            \item Verify your internet connection
            \item Check if Lodestone is undergoing maintenance
        \end{itemize}
    \item If GitHub Pages deployment fails:
        \begin{itemize}
            \item Ensure the repository is public
            \item Check if the branch and folder settings are correct
        \end{itemize}
\end{itemize}

\section{Development Notes}
\begin{itemize}
    \item The application uses React for the UI
    \item Data is fetched directly from Lodestone
    \item Styling uses CSS modules for component isolation
    \item Webpack handles bundling and development server
\end{itemize}

\end{document} 